\section{Finite Signed Graphs and Binary Extension Constraints}
\label{fin:chapter}

The models in this section are finite mathematical comparators.  The
graph, real or binary variables, variational functional, probability
measure, and any quantum representation are specified separately.
Absence of a time parameter from a functional is not a derivation of
physical activity without time.  In particular, a variational minimum
does not itself choose a sampling law or an observation instrument.

\subsection{A graph double well and its exact threshold}
\label{fin:graph-action}

For a finite undirected simple graph \(G=(V,E)\), let \(A_G,D_G\)
be its adjacency and degree matrices and \(Q_G=D_G+A_G\) its signless
Laplacian.  For real amplitudes \(a_i\), define
\[
 S(a)=\frac{\kappa}{2}\sum_{\{i,j\}\in E}(a_i+a_j)^2
       +\frac{\lambda}{4}\sum_i(a_i^2-m^2)^2,
 \qquad \kappa,\lambda,m>0 .
\]
The signless-Laplacian quadratic form is standard graph theory
\cite{fin:signless}.

\begin{theorem}[Threshold, bipartiteness, and balance]\label{fin:threshold}
Let \(q_{\min}\) be the least eigenvalue of \(Q_G\).
The origin is the unique global minimum of \(S\) if
\(\kappa q_{\min}\ge\lambda m^2\).
If \(\kappa q_{\min}<\lambda m^2\), global nonzero minimizers exist and
the origin is not a global minimum.
For connected \(G\), a zero-action configuration exists exactly when
\(G\) is bipartite.  In that case its only zero-action configurations
are the two alternating assignments of magnitude \(m\).
They have zero signed total exactly when the two bipartition classes
have equal size.
\end{theorem}
\begin{proof}
Expanding the squares gives
\[
 S(a)-S(0)=\frac12a^T(\kappa Q_G-\lambda m^2 I)a
                    +\frac{\lambda}{4}\sum_i a_i^4 .
\]
When the quadratic form is nonnegative, the quartic term is strictly
positive for every nonzero \(a\), even at equality of the threshold.
When the least quadratic coefficient is negative, substitute \(a=tv\)
for a corresponding unit eigenvector.  The negative order-\(t^2\)
term dominates the positive order-\(t^4\) term for sufficiently small
nonzero \(t\).  The functional is coercive, so it attains its global
minimum, which is then nonzero.
Every term of the original functional is nonnegative.  Zero action
requires \(|a_i|=m\) and \(a_i=-a_j\) on every edge.
On a connected graph these relations propagate a unique sign choice
up to global reversal and are consistent precisely when there is no
odd cycle, equivalently when the graph is bipartite.
For bipartition \(U\sqcup W\), the signed sums are
\(\pm m(|U|-|W|)\).
\end{proof}

On the four-cycle the minima are \(\pm(m,-m,m,-m)\).
Their Hessian is \(\kappa Q_G+2\lambda m^2I\), with eigenvalues
\[
 2\lambda m^2,\quad
 2\kappa+2\lambda m^2\ \text{(twice)},\quad
 4\kappa+2\lambda m^2 .
\]
The signed sum is zero, the mean square is \(m^2\), and
\(\sum_{\{i,j\}\in E}(a_i-a_j)^2=16m^2\).
These are positive relational variation and a strict variational
minimum with zero action.  They do not describe evolution under a
physical equation of motion.  On \(K_{1,3}\), the zero-action
configurations instead have total signed amplitudes \(\pm2m\).
An equally weighted ensemble of them has mean zero, which is a
different statement from configurationwise balance.

If the amplitudes are additionally restricted to \(a_i=m\sigma_i\),
\(\sigma_i\in\{-1,+1\}\), the onsite term vanishes and
\[
 S(\sigma)=2\kappa m^2 f_G(\sigma),\qquad
 \min_\sigma S=2\kappa m^2(|E|-\operatorname{MaxCut}(G)),
\]
where \(f_G\) counts equal-sign edges.  This discrete restriction and
its Ising/MaxCut problem are not the unrestricted continuous
minimization problem.

\subsection{Action-count statistics and graph identification}
\label{fin:partitions}

An additional Gibbs stipulation gives
\[
 Z_G(z)=\sum_{\sigma\in\{\pm1\}^V}z^{f_G(\sigma)},\qquad
 p_G(\sigma)=z^{f_G(\sigma)}/Z_G(z),\qquad z>0.
\]
No probability law has been inferred from the extrema alone.

\begin{proposition}[Trees, cycles, and the even-subgraph formula]
\label{fin:partition-theorem}
Every \(n\)-vertex tree has \(Z_G(z)=2(1+z)^{n-1}\).
For a cycle of length \(n\),
\[
 Z_{C_n}(z)=(1+z)^n+(z-1)^n .
\]
For any finite graph,
\[
 Z_G(z)=2^{|V|}\left(\frac{1+z}{2}\right)^{|E|}
 \sum_{\substack{F\subseteq E\\\deg_F(v)\text{ even for all }v}}
       \left(\frac{z-1}{z+1}\right)^{|F|}.
\]
\end{proposition}
\begin{proof}
On a rooted tree, choose the root sign and, independently as
combinatorial coordinates, whether each edge joins equal signs.
These choices bijectively determine the entire configuration.
Summing their weights gives \(2(1+z)^{n-1}\).
On a cycle the transfer matrix
\(\left(\begin{smallmatrix}z&1\\1&z\end{smallmatrix}\right)\)
has eigenvalues \(1+z,z-1\); its \(n\)th-power trace counts all closed
sign sequences and gives the stated formula.
Finally, for each edge,
\[
 z^{\mathbf1_{\{\sigma_i=\sigma_j\}}}
 =\frac{1+z}{2}\left(1+\frac{z-1}{z+1}\sigma_i\sigma_j\right).
\]
Expand the product over edges.  Summing over each vertex sign kills
every term in which its degree in the selected edge set is odd;
all remaining terms contribute \(2^{|V|}\).
\end{proof}

Summing out a pendant vertex contributes \(1+z\), regardless of where
it was attached.  In particular, a connected unicyclic graph with
cycle length \(\ell\) has
\((1+z)^{n-\ell}Z_{C_\ell}(z)\).
Thus action-count statistics can lose graph shape.  Such
partition-function equivalence is an established graph-identification
issue \cite{fin:ising}.

\subsubsection{A stipulated quantum observation of the action count}

To make one observation map explicit, take an orthonormal configuration
register and a two-dimensional sign register, and stipulate
\[
 |\Psi_G\rangle=\sum_\sigma\sqrt{p_G(\sigma)}\,|\sigma\rangle
          |\chi_{f_G(\sigma)}\rangle,\qquad
 |\chi_f\rangle=
 \frac{e^{i\theta f}|+\rangle+e^{-i\theta f}|-\rangle}{\sqrt2}.
\]
The Hilbert spaces, phase rule, tensor product and Born measurements are
additional inputs.

\begin{proposition}[Reduced sign law and conditional mixedness]
\label{fin:graph-readout}
The reduced sign state and two-copy SWAP-test probability are
\[
 \rho_G=\frac12\begin{pmatrix}1&C_G\\\overline C_G&1\end{pmatrix},
 \qquad C_G=\frac{Z_G(ze^{2i\theta})}{Z_G(z)},\qquad
 p_{\rm SWAP+}=\frac{3+|C_G|^2}{4}.
\]
Graphs with identical \(Z_G\) have identical sign-register laws for
all common parameters and all measurements on independently prepared
copies.  The pure conditional-sign ansatz is not in general closed
under eliminating graph vertices by ordinary partial trace.
\end{proposition}
\begin{proof}
Tracing the orthonormal configuration labels gives
\(\rho_G=\sum_\sigma p_G(\sigma)|\chi_f\rangle\langle\chi_f|\).
Its off-diagonal entry is half the characteristic function
\(E e^{2i\theta f}\), proving the ratio formula.
The SWAP identity is
\(p_{\rm SWAP+}=(1+\operatorname{Tr}\rho_G^2)/2\), and
\(\operatorname{Tr}\rho_G^2=(1+|C_G|^2)/2\).
Equality of the reduced state gives equality of all its independently
prepared tensor powers.

For the failure of closure, eliminate vertex three of a triangle and
fix either fully specified equal-sign configuration of the retained
vertices.  The eliminated sign has action counts \(f=1,3\) with
relative weights \(z,z^3\).  The conditional visibility is
\[
 C_{\rm eq}=\frac{e^{2i\theta}+z^2e^{6i\theta}}{1+z^2},\qquad
 |C_{\rm eq}|^2=
 \frac{1+z^4+2z^2\cos(4\theta)}{(1+z^2)^2}.
\]
For \(z>0\) and \(e^{4i\theta}\ne1\), this is strictly less than one.
The resulting conditional sign state is mixed, whereas the original
ansatz assigns a pure sign state to every complete graph configuration.
Retuning the smaller graph's parameters cannot remove that difference
while retaining this pure conditional ansatz.
\end{proof}

At \(z=1/2,\theta=\pi/8\), the equal retained-sign event has probability
\(5/13\); its conditional state has \(|C_{\rm eq}|^2=17/25\) and
purity \(21/25\).  Its distance from the nearest pure qubit state is
\((5-\sqrt{17})/10\), since the Bloch-vector length is \(\sqrt{17}/5\).
For the complete triangle and square sign states, respectively,
\[
 |C_{C_3}|^2=145/169,\quad |C_{C_4}|^2=801/1681,\qquad
 p_{C_3}=163/169,\quad p_{C_4}=1461/1681.
\]
An enlarged family of arbitrary joint density matrices is closed under
partial trace; this is a valid positive comparison, but it replaces
the restricted pure conditional ansatz with ordinary mixed-state
quantum theory.

The reduced state already has a conventional random-phase preparation:
\[
 \sum_fP(f)e^{i\theta fZ}|+\!x\rangle\langle+\!x|e^{-i\theta fZ}.
\]
It is not a classical simulation of all quantum experiments.
Writing \(F,F'\) for independent action counts also gives
\[
 1-|C_G|^2=E[1-\cos(2\theta(F-F'))]
 \le4\theta^2\operatorname{Var}(F).
\]
Thus \(1-p_{\rm SWAP+}\le\theta^2\operatorname{Var}(F)\);
the leading small-phase term is
\(\theta^2\operatorname{Var}(F)\).
At \(\theta=0\), or as \(z\to0^+\) for a fixed finite graph, the
SWAP probability tends to one.  The nonzero finite-parameter contrast
does not persist uniformly in those limits.

\begin{example}[Graph aliases with different squared magnetization]
\label{fin:graph-aliases}
Let \(M=\sum_i\sigma_i\).
On vertices \(0,\ldots,4\), take edge sets
\[
 A_5=\{01,12,20,03,04\},\qquad
 B_5=\{01,12,20,03,34\}.
\]
Their degree sequences differ, so they are nonisomorphic, but
both have \(Z_5=2z(1+z)^2(3+z^2)\).
On vertices \(0,\ldots,5\), take
\[
 A_6=\{01,05,12,15,23,24,34,45\},\qquad
 B_6=\{01,04,05,12,15,23,25,34\}.
\]
Both have degree sequence \((3,3,3,3,2,2)\).
Their two degree-two vertices are nonadjacent in \(A_6\) and adjacent
in \(B_6\), proving nonisomorphism.  Both have
\[
 Z_6=2z^2(1+z)^2(1+z^2)(z^2-2z+5).
\]
Here an edge \(ij\) denotes the unordered pair \(\{i,j\}\).
\end{example}

For a complete finite check of the last claims, let
\(W_G(u)=\sum_{\sigma:M^2=u}z^{f_G(\sigma)}\).
Fixing the sign at vertex zero and doubling enumerates all
configurations.  The resulting polynomials for the two \(A\) graphs are
\[
\begin{array}{c|l}
u&W_{A_5}(u)\\ \hline
1&4z^3+12z^2+4z\\
9&4z^4+4z^3+2z\\
25&2z^5
\end{array}
\qquad
\begin{array}{c|l}
u&W_{A_6}(u)\\ \hline
0&2z^6+4z^4+8z^3+6z^2\\
4&8z^5+10z^4+8z^3+4z^2\\
16&4z^6+8z^5\\
36&2z^8 .
\end{array}
\]
For the five-vertex pair the only differences \(W_{A_5}-W_{B_5}\)
are \(+2z(1+z)(1-z)^2\) at \(u=9\) and its negative at \(u=1\).
For the six-vertex pair they are \(+2z^2(1+z)(1-z)^3\) at \(u=4\)
and its negative at \(u=0\).
These identities specify all the \(B\) polynomials as well and prove
the equal partition functions by summation.  Dividing by the common
normalization gives, for \(0<z\le1\),
\[
 \operatorname{TV}_5=
 \frac{(1-z)^2}{(1+z)(3+z^2)},\qquad
 \operatorname{TV}_6=
 \frac{(1-z)^3}{(1+z)(1+z^2)(z^2-2z+5)}.
\]
At \(z=1/2\), these are \(2/39\) and \(4/255\).
The differences of \(E M^2\) are respectively eight and four times
these distances.  The six-vertex distance decays cubically as \(z\to1\),
with leading term \((1-z)^3/16\).
Observing \(M^2\) uses configuration-sensitive information absent from
the reduced sign register.  These are exact distributional contrasts,
not measured effects or instrumental precision claims.

\subsubsection{Complete labelled sign observations recover adjacency}
\label{fin:labelled-sign-recovery}

The graph aliases above concern a restricted observation.  There is a
positive identification result when the entire labelled sign law
\(p_G(\sigma)=z^{f_G(\sigma)}/Z_G(z)\) is available, with known
\(0<z<1\) and a finite simple graph on a fixed labelled vertex set.
Fix distinct vertices \(i,j\) and any assignment \(\tau\) of all other
signs.  Write \(p_{ab}=p_G(\sigma_i=a,\sigma_j=b,\tau)\) for
\(a,b\in\{+,-\}\).  Strict positivity makes the ratio well defined, and
\[
 \frac{p_{++}p_{--}}{p_{+-}p_{-+}}
       =z^{\,2A_{ij}},\qquad
 A_{ij}=\frac{\log(p_{++}p_{--}/p_{+-}p_{-+})}{2\log z}.
\]
Indeed, an edge not joining \(i\) and \(j\) contributes equally to the
two numerator and two denominator exponents: edges to a fixed outside
sign have one equal and one unequal endpoint assignment in either
product.  The edge \(ij\), if present, contributes two in the numerator
and zero in the denominator.  Normalization cancels.  Thus the ratio
is \(z^2\) on an edge and one otherwise, recovering every labelled
adjacency entry.

At \(z=1\) the law is uniform for every graph, so the contrast vanishes.
The zero-temperature endpoint can have support zeros and does not obey
this positivity argument.  Unlabelled observations or the reduced
action-count register do not supply these four probabilities.  Exact
population recovery therefore preserves the graph-alias counterexamples
and makes no finite-sample or physical-access guarantee.

\subsection{Opposite-sign adjacency and an operator-dependent spectrum}
\label{fin:spectra}

\begin{proposition}[Complete-bipartite spectra]\label{fin:bipartite-spectrum}
The rule \(i\sim j\) exactly when \(\sigma_i\ne\sigma_j\) produces
an empty graph or a complete bipartite graph \(K_{p,q}\).
For \(p,q>0\), \(N=p+q\), the symmetric normalized Laplacian has spectrum
\[
 \{0,\ 1^{(N-2)},\ 2\},
\]
whereas the unnormalized Laplacian has eigenvalues
\[
 0,\quad q\ \text{with multiplicity }p-1,\quad
 p\ \text{with multiplicity }q-1,\quad N.
\]
\end{proposition}
\begin{proof}
The two sign classes determine all edges, giving \(K_{p,q}\).
Its normalized adjacency matrix has off-diagonal blocks
\(J_{p,q}/\sqrt{pq}\) and its transpose.  On vectors constant within
each part its eigenvalues are \(+1,-1\); on the within-part
zero-sum spaces it vanishes.  Subtracting from identity proves the
normalized spectrum.
For the unnormalized Laplacian, a vector supported on the first part
and summing to zero has eigenvalue \(q\), and analogously the other
part gives \(p\).  The remaining two-dimensional constant-on-parts
space has eigenvalues \(0,N\).
\end{proof}

Thus \(K_{1,3}\) and \(K_{2,2}\) are normalized-Laplacian isospectral,
while their unnormalized spectra are \((0,1,1,4)\) and \((0,2,2,4)\).
The normalized heat trace and its formal dimension diagnostic are
\[
 K_N(s)=\frac{1+(N-2)e^{-s}+e^{-2s}}N,\qquad
 d_{\rm spec}(s)=
 \frac{2s[(N-2)e^{-s}+2e^{-2s}]}
      {1+(N-2)e^{-s}+e^{-2s}}.
\]
At each fixed \(s>0\), \(d_{\rm spec}(s)\to2s\) as \(N\to\infty\).
Choosing \(s=1\) gives the number two while the ordinary graph diameter
stays at most two.  There is no scale-independent plateau in this
fixed-\(s\) limit.  A chosen graph operator and one diffusion-scale
number therefore do not establish a physical spatial dimension or
a local continuum.

\subsection{One four-sign law and its exact extensions}
\label{fin:four-law}

For four signs, put \(M_4=\sum_i\sigma_i\) and
\[
 P_4(\sigma)=\frac{2^{-M_4^2/4}}{81/8}.
\]
Let \(K_N\) count the plus signs among \(N\) coordinates.  Grouping
the sixteen configurations by \(K_4\) gives
\begin{equation}\label{fin:target}
 P_4(K_4=k)=
 \left(\frac1{162},\frac{16}{81},\frac{16}{27},
                    \frac{16}{81},\frac1{162}\right)_k .
\end{equation}
The law is exchangeable and invariant under global sign reversal.
Its signed mean is zero and \(E M_4^2=16/9\).
Since \(M_4^2=4+2\sum_{i<j}\sigma_i\sigma_j\),
every distinct pair has correlation \(-5/27\).
Finite extendibility is a conventional probability problem
\cite{fin:extendibility}; its requirements should be stated for
complete marginals, not inferred from a few moments.

\begin{theorem}[Unique five- and six-coordinate extensions]
\label{fin:rigidity}
If every four-coordinate marginal of a law on \(N=5\) or \(6\)
labelled signs equals the complete law \(P_4\), then the larger law
is unique, even without assuming exchangeability.
For \(N=5\), each configuration with \(K_5=2\) or \(3\) has probability
\(4/81\); the all-minus and all-plus configurations each have probability
\(1/162\); all others have probability zero.
For \(N=6\), each configuration with \(K_6=3\) has probability \(4/81\),
the two all-equal configurations have probability \(1/162\), and all
others have probability zero.
\end{theorem}
\begin{proof}
Let \(p_j\) be the count law of a uniformly selected four-coordinate
subset and \(q_j=P(K_N=j)\).
Conditional on \(K_N=k\), this sampled count has probability
\(\binom kj\binom{N-k}{4-j}/\binom N4\), regardless of exchangeability.
For \(N=5\), the nonnegative integer-count polynomial
\[
 Q_5(k)=k(5-k)(k-2)(k-3)
\]
equals \(8\) at \(k=1,4\) and zero otherwise.  Direct substitution in
the hypergeometric probabilities gives
\[
 E Q_5(K_5)=8(q_1+q_4)
          =10(p_1+p_3)-\frac{20}{3}p_2 .
\]
The right side is zero for \eqref{fin:target}, so the support is
\(\{0,2,3,5\}\).
A specified four-coordinate pattern with one plus and three minuses
has probability \(4/81\).  Its missing sign must be plus because
total count one is forbidden.  This fixes each individual
two-plus configuration at \(4/81\).  Global sign reversal of the
argument fixes the three-plus configurations; no symmetry assumption
on the candidate law is needed because the prescribed marginal itself
has the reversed pattern probabilities.  An all-plus four-pattern
can extend only to the all-plus five-pattern, fixing its mass at
\(1/162\); likewise for all minus.

For \(N=6\), use
\[
 Q_6(k)=k(6-k)(k-3)^2,\qquad
 E Q_6(K_6)=20(q_1+q_5)+8(q_2+q_4)
          =30(p_1+p_3)-20p_2 .
\]
It is nonnegative and its expectation is zero, forcing support
\(\{0,3,6\}\).
A four-coordinate pattern with three pluses and one minus then forces
both omitted signs to be minus.  This fixes every three-plus
configuration at \(4/81\), and the extreme patterns fix the remaining
two masses at \(1/162\).
For existence in either size, the listed probabilities sum to one.
For a specified four-pattern with \(k=0,1,2,3,4\) pluses, summing over
the omitted coordinates gives, respectively,
\((1/162,4/81,8/81,4/81,1/162)\) per pattern, exactly \(P_4\).
\end{proof}

The corresponding count masses are
\[
 \begin{array}{c|l}
 N=5&(q_0,q_2,q_3,q_5)=(1/162,40/81,40/81,1/162)\\
 N=6&(q_0,q_3,q_6)=(1/162,80/81,1/162).
 \end{array}
\]
This conclusion requires all four-coordinate marginals.
A single designated marginal always has an extension by
adjoining an independent sign, and no uniqueness assertion applies to
that different problem.
The displayed exact extensions require zero probabilities.
Everywhere finite Gibbs actions with counting measure give strictly
positive probabilities, so cannot realize these exact extensions.

\begin{corollary}[No seven-coordinate exact extension]\label{fin:no-seven}
No law on seven labelled signs has every four-coordinate marginal
equal to \(P_4\).
\end{corollary}
\begin{proof}
Every pair would have correlation \(-5/27\), hence
\[
 E M_7^2=7+42(-5/27)=-7/9,
\]
contradicting nonnegativity.
\end{proof}

\subsection{Sharp finite and infinite approximation distances}
\label{fin:distances}

\begin{theorem}[Odd-even finite-extension staircase]\label{fin:staircase}
Let \(N=2m-1\) or \(2m\), \(m\ge4\).
The minimum total variation from \(P_4\) to the four-coordinate
marginal of an exchangeable \(N\)-sign law is
\[
 d_N=\frac{16}{27}
 -\frac{3m(m-1)}{2(2m-1)(2m-3)}.
\]
In particular \(d_7=d_8=74/945\).
The corresponding minimum over infinitely exchangeable laws is
\(47/216\).
\end{theorem}
\begin{proof}
Conditional on a population total \(k\), the balance event \(B=\{K_4=2\}\)
has probability
\[
 h_N(k)=\frac{\binom k2\binom{N-k}2}{\binom N4}.
\]
For interior \(k\), the ratio \(h_N(k+1)/h_N(k)\) is
\((k+1)(N-k-2)/[(k-1)(N-k)]\); it is at least one precisely up to
the central count.  Thus its maximum is attained at the central
count or counts and equals
\[
 h_{\max}=\frac{3m(m-1)}{2(2m-1)(2m-3)}.
\]
Every mixture over counts obeys the same bound.
Since \(P_4(B)=16/27\), the balance-event deficit is a total variation
lower bound.

For even \(N=2m\), choose exactly \(m\) pluses uniformly over positions.
For odd \(N=2m-1\), mix the two central counts \(m-1,m\) equally and
uniformly within each count.  Both have the same four-count law:
\[
 q_0=q_4=\frac{(m-2)(m-3)}
                    {4(2m-1)(2m-3)},\qquad
 q_1=q_3=\frac{m(m-2)}{(2m-1)(2m-3)},\qquad q_2=h_{\max}.
\]
For \(m\ge4\), the outer probabilities are at least \(1/70>1/162\),
the next pair at least \(8/35>16/81\), and the center is below \(16/27\).
For example, the first two lower comparisons reduce respectively to
\((m-4)(9m-17)\ge0\) and \((m-4)(m+2)\ge0\).
Only the central category is deficient, so total variation equals
the central deficit.  Exchangeable laws are uniform within each count
class, so their full-pattern total variation equals their count-law
total variation.

For an infinite exchangeable law, restriction to every \(N\) yields
the same four-law, so the finite balance bounds imply
\(Q(B)\le\lim_Nh_{\max}=3/8\).
Consequently the distance is at least \(16/27-3/8=47/216\).
Iid fair signs attain it: their count law
\((1,4,6,4,1)/16\) has all four noncentral categories above the target.
\end{proof}

The next plateaus are \(d_9=d_{10}=22/189\) and
\(d_{11}=d_{12}=41/297\), increasing to \(47/216\).
For an arbitrary, possibly nonexchangeable \(N\)-sign law, uniformly
sampling four indices still obeys the same balance-event bound.
At least one four-coordinate marginal must therefore be at least
\(d_N\) from the target, by data processing and convexity.
These are distances between ideal distributions.  They are not
experimental detector tolerances or physical population limits.

\subsection{A sharp bounded-action approximation}
\label{fin:bounded-action}

Let \(R\ge1\).  Consider all strictly positive probability laws \(p\)
on five signs satisfying
\(\max_\sigma p(\sigma)/\min_\sigma p(\sigma)\le R\).
For a Gibbs law \(p\propto e^{-s(\sigma)}\) relative to counting measure,
this is exactly an action-range bound
\(\max s-\min s\le\log R\).
Let \(d_*(R)\) be the minimum, over such laws, of the largest total
variation distance between any of their five complete four-coordinate
marginals and \(P_4\).

\begin{theorem}[Exact action-range tradeoff]\label{fin:range-theorem}
The minimum is attained and equals
\[
 d_*(R)=
 \begin{cases}
 \displaystyle\frac{16}{27}-\frac{3R}{3+5R},
                       &1\le R\le78/5,\\[6pt]
 \displaystyle\frac{32}{27(5R+2)},&R\ge78/5 .
 \end{cases}
\]
\end{theorem}
\begin{proof}
Average any admissible law over coordinate permutations and global
sign reversal.  The constraints \(p(\sigma)\le R p(\tau)\) are linear
and invariant under these transformations, so the ratio bound
survives.  Convexity of total variation and symmetry of the target
show that the largest marginal error cannot increase.
Thus it suffices to consider per-configuration probabilities \(a,b,c\)
for count classes \(\{0,5\},\{1,4\},\{2,3\}\), respectively.
They obey
\[
 2a+10b+20c=1,\qquad
 q=(a+b,\ 4(b+c),\ 12c,\ 4(b+c),\ a+b).
\]

First, \(a,b\ge c/R\) implies
\(c\le R/[4(5R+3)]\).  The central event therefore gives
\[
 \operatorname{TV}(P_4,q)\ge16/27-3R/(5R+3).
\]
Set
\[
 a=b=\frac1{4(5R+3)},\qquad c=\frac{R}{4(5R+3)}.
\]
These weights are positive, normalized and have ratio \(R\).
For \(1\le R\le78/5\), \(q_0=q_4\ge1/162\) and
\(q_1=q_3>16/81\), while \(q_2<16/27\).
Hence the total variation is exactly the central deficit, attaining
the first formula.

For the second lower bound, put \(v=3R/(5R+2)\in[0,1]\) and use
the test function \(f(2)=1\), \(f(0)=f(4)=v\), \(f(1)=f(3)=0\).
After eliminating \(a\) with normalization,
\[
 E_qf=12c+2v(a+b)
     =v+\frac{24(c-Rb)}{5R+2}\le v.
\]
Since \(E_{P_4}f=16/27+v/81\), the bounded-function characterization
of total variation gives
\[
 \operatorname{TV}(P_4,q)\ge16/27+v/81-v
                          =\frac{32}{27(5R+2)}.
\]
For \(R\ge78/5\), choose
\[
 b=\frac{20}{81(5R+2)},\qquad
 c=Rb,\qquad
 a=\frac1{162}-b=\frac{5R-38}{162(5R+2)}.
\]
They are positive and normalized.  Here \(b\le a\le c\), so their
maximum-to-minimum ratio is \(R\).
The outer marginal masses exactly match \(1/162\);
each near-central mass exceeds its target by
\(16/[27(5R+2)]\), while the central deficit is
\(32/[27(5R+2)]\).  The lower bound is attained.
At \(R=78/5\) the constructions and formulas agree.
\end{proof}

Thus the error is positive for every finite action range and decays
as \((32/135)R^{-1}+O(R^{-2})\) for large \(R\).
This concerns all four-subset marginals simultaneously.  An exact
extension of only one designated marginal can instead be formed by
adjoining an independent sign, with finite positive weights.
The probability-ratio optimization does not postulate a physical
action law or calibrate an instrument.

\subsection{Quadratic weights are not closed under deletion}
\label{fin:quadratic}

\begin{proposition}[A five-to-four closure obstruction]
\label{fin:quadratic-closure}
For \(b>0\), the size-dependent quadratic family
\[
 P_N(\sigma)\ \propto\ \exp[-b M_N^2/(2N)]
\]
is not projectively consistent under ordinary coordinate deletion.
Even allowing an adjusted quadratic coefficient at size four does not
make the five-to-four marginal a pure quadratic weight.
\end{proposition}
\begin{proof}
For retained magnetization \(M\in\{0,\pm2,\pm4\}\), summing the fifth
sign gives
\[
 w_b(M)=2e^{-b(M^2+1)/10}\cosh(bM/5).
\]
Every weight \(w(M)=C e^{-\alpha M^2}\) necessarily satisfies
\(w(4)w(0)^3/w(2)^4=1\), because both its constant and quadratic
exponents cancel.  Here the ratio is
\[
 \frac{\cosh(4b/5)}{\cosh^4(2b/5)}
 =1-\tanh^4(2b/5)<1.
\]
To obtain the identity, set \(x=2b/5\), use
\(\cosh(2x)=2\cosh^2x-1\), and divide by \(\cosh^4x\).
\end{proof}

The multiplicative defect is \(16b^4/625+O(b^6)\), so a lower-order
expansion can miss it.  At \(b=0\) the iid fair family is exactly
projective.  There is also a small-size exception: after reducing a
three-sign law to two signs, only two retained values of \(M^2\) remain, so a
normalization and a retuned quadratic coefficient can fit both
positive weights.  Three distinct \(M^2\) values in the five-to-four
experiment impose the extra compatibility equation that fails.
The valid bare probability extensions of Theorem~\ref{fin:rigidity}
are therefore not generated by simply reusing the positive quadratic
action.  Their existence and this particular action's failure are
compatible facts.
